
\section{Crossings (take 2)}

This is a work-in-progress to give revised proofs of the correct
conditions for crossings to be essential.

\begin{definition}
  Recall some facts about connecting curves with RIII moves to make
  crossings in bijection and make precise the topological notion of
  \emph{essential crossing}. Define (topological) smoothing.
\end{definition}

\begin{definition}
  Non-trivial group elements $a,b\in \pi_1$ are \emph{common powers}
  if there are non-zero $k,\ell$ so that $a^k = b^\ell$. Since $\pi_1$
  has no torsion, we may assume that $k,\ell$ are relatively prime and
  then we see this is equivalent to the existence of a $c$ so that
  $a = c^\ell$ and $b = c^k$.
\end{definition}

Define embedding of 4-valent graph $[a * b]$. Homotopic under
simultaneous conjugation of $a,b$. \emph{Taut} still means minimal
intersections.

\begin{theorem}
  Let $S$ be a surface with basepoint~$*$ and $a,b \in \pi_1(S,*)$ be
  any non-trivial elements that are not common powers. Then in the
  following three trialities the same possibility holds in each case.
  \begin{enumerate}
  \item (Reduction) Exactly one of the following multi-curve reduces to the other two:\\
    \begin{enumerate*}[itemjoin=\qquad]
    \item $[ab]$;
    \item $[aB]$; or
    \item $[a][b]$.
    \end{enumerate*}
  \item (Geometric) For every taut embedding of $a * b$, around the
    4-valent vertex we see in cyclic
    order (up to orientation reversal):\\
    \begin{enumerate*}
    \item $a$ in, $a$ out, $b$ out, $b$ in;
    \item $a$ in, $a$ out, $b$ in, $b$ out; or
    \item $a$ in, $b$ in, $a$ out, $b$ out
    \end{enumerate*}
  \item (Algebraic) The geodesics $\ora{a}$ and $\ora{b}$ are:\\
    \begin{enumerate*}[itemjoin=\qquad]
    \item parallel;
    \item anti-parallel; or
    \item cross.
    \end{enumerate*}
  \end{enumerate}
\end{theorem}

\begin{proof}
  Note that the algebraic triality is clearly mutually exclusive and
  exhaustive (since the endpoints of $\ora{a}$ and $\ora{b}$ are
  distinct since they are not common powers). For the geometric triality,
  for any given embedding of $a * b$, one of the three given
  possibilities holds, but it is not a priori clear that the same
  possibility holds for all taut embeddings. For the reduction triality,
  the possibilities are mutually exclusive (we can't have circular
  reductions, since at each reduction length with respect any
  Riemannian metric strictly decreases) but it is not a priori clear
  that any of the three possibilities holds.
  
  We start with the geometric triality. Consider any taut embedding
  $\Gamma$ of $[a * b]$ in one of the cases (a), (b), or~(c). Consider
  the (concrete) multi-curve~$C$ given by replacing the
  $4$-valent vertex with a crossing. If we pick orientations on $C$
  correctly, this is an embedding of $[ab]$, $[aB]$, or $[a][b]$,
  respectively, with a distinguished crossing~$x$.

  We claim that $C$ is taut. Otherwise, there is some reducing isotopy
  of $C$, using only RIII moves and some RI/RII reduction. Any RIII move on~$C$
  can be copied with an isotopy of~$\Gamma$, as can an RI or RII reduction of
  $C$ that does not involve~$x$. An RI
  reduction at~$x$ is only possible if $a$ or $b$ is the identity, and
  an RII reduction involving~$x$ also reduces $\Gamma$, though
  reducing the number of intersections of $\Gamma$ by one rather than
  by two. In any case this contradicts tautness of~$\Gamma$ or
  non-triviality of $a,b$.

  We have thus shown that one of $[ab]$, $[aB]$, or $[a][b]$ has a
  taut representative $C$ with an essential crossing~$x$ that can be
  smoothed to give the other two multi-curves. We cannot have a
  circular chain of essential smoothings (since every essential
  smoothing strictly decreases self-intersection number) so it follows
  that the cases are exclusive and thus that \emph{every} taut
  representative of $[a * b]$ has the same cyclic order around the
  4-valent vertex, proving equivalence of the first two trialities.

  For the remaining equivalence, suppose first we
  are in case~(c) of the geometric triality. For the fixed hyperbolic
  structure, take $C$ to be the geodesic representative of the free
  homotopy class $[a][b]$, intersecting at the distinguished point
  $x \in \Sigma$, and let $\ora{a}$ and $\ora{b}$ be the respective
  lifts passing through some lift $\wt{x}$ of~$x$. Then $\ora{a}$ and
  $\ora{b}$ intersect only once and thus cross at infinity, so we are
  also in case~(c) of the algebraic triality.
  
  \begin{itemize}
  \item In the parallel case (a), again take the geodesic
    representative, but now we need to argue that after the smoothing
    there is still no monogon. This is exactly claim 4 from my notes
    \texttt{2024-06-03 \#3}.
  \end{itemize}
\end{proof}

%%% Local Variables:
%%% mode: latex
%%% TeX-master: "Intersections"
%%% End:
